% Options for packages loaded elsewhere
% Options for packages loaded elsewhere
\PassOptionsToPackage{unicode}{hyperref}
\PassOptionsToPackage{hyphens}{url}
\PassOptionsToPackage{dvipsnames,svgnames,x11names}{xcolor}
%
\documentclass[
  letterpaper,
  DIV=11,
  numbers=noendperiod]{scrreprt}
\usepackage{xcolor}
\usepackage{amsmath,amssymb}
\setcounter{secnumdepth}{5}
\usepackage{iftex}
\ifPDFTeX
  \usepackage[T1]{fontenc}
  \usepackage[utf8]{inputenc}
  \usepackage{textcomp} % provide euro and other symbols
\else % if luatex or xetex
  \usepackage{unicode-math} % this also loads fontspec
  \defaultfontfeatures{Scale=MatchLowercase}
  \defaultfontfeatures[\rmfamily]{Ligatures=TeX,Scale=1}
\fi
\usepackage{lmodern}
\ifPDFTeX\else
  % xetex/luatex font selection
\fi
% Use upquote if available, for straight quotes in verbatim environments
\IfFileExists{upquote.sty}{\usepackage{upquote}}{}
\IfFileExists{microtype.sty}{% use microtype if available
  \usepackage[]{microtype}
  \UseMicrotypeSet[protrusion]{basicmath} % disable protrusion for tt fonts
}{}
\makeatletter
\@ifundefined{KOMAClassName}{% if non-KOMA class
  \IfFileExists{parskip.sty}{%
    \usepackage{parskip}
  }{% else
    \setlength{\parindent}{0pt}
    \setlength{\parskip}{6pt plus 2pt minus 1pt}}
}{% if KOMA class
  \KOMAoptions{parskip=half}}
\makeatother
% Make \paragraph and \subparagraph free-standing
\makeatletter
\ifx\paragraph\undefined\else
  \let\oldparagraph\paragraph
  \renewcommand{\paragraph}{
    \@ifstar
      \xxxParagraphStar
      \xxxParagraphNoStar
  }
  \newcommand{\xxxParagraphStar}[1]{\oldparagraph*{#1}\mbox{}}
  \newcommand{\xxxParagraphNoStar}[1]{\oldparagraph{#1}\mbox{}}
\fi
\ifx\subparagraph\undefined\else
  \let\oldsubparagraph\subparagraph
  \renewcommand{\subparagraph}{
    \@ifstar
      \xxxSubParagraphStar
      \xxxSubParagraphNoStar
  }
  \newcommand{\xxxSubParagraphStar}[1]{\oldsubparagraph*{#1}\mbox{}}
  \newcommand{\xxxSubParagraphNoStar}[1]{\oldsubparagraph{#1}\mbox{}}
\fi
\makeatother


\usepackage{longtable,booktabs,array}
\newcounter{none} % for unnumbered tables
\usepackage{calc} % for calculating minipage widths
% Correct order of tables after \paragraph or \subparagraph
\usepackage{etoolbox}
\makeatletter
\patchcmd\longtable{\par}{\if@noskipsec\mbox{}\fi\par}{}{}
\makeatother
% Allow footnotes in longtable head/foot
\IfFileExists{footnotehyper.sty}{\usepackage{footnotehyper}}{\usepackage{footnote}}
\makesavenoteenv{longtable}
\usepackage{graphicx}
\makeatletter
\newsavebox\pandoc@box
\newcommand*\pandocbounded[1]{% scales image to fit in text height/width
  \sbox\pandoc@box{#1}%
  \Gscale@div\@tempa{\textheight}{\dimexpr\ht\pandoc@box+\dp\pandoc@box\relax}%
  \Gscale@div\@tempb{\linewidth}{\wd\pandoc@box}%
  \ifdim\@tempb\p@<\@tempa\p@\let\@tempa\@tempb\fi% select the smaller of both
  \ifdim\@tempa\p@<\p@\scalebox{\@tempa}{\usebox\pandoc@box}%
  \else\usebox{\pandoc@box}%
  \fi%
}
% Set default figure placement to htbp
\def\fps@figure{htbp}
\makeatother





\setlength{\emergencystretch}{3em} % prevent overfull lines

\providecommand{\tightlist}{%
  \setlength{\itemsep}{0pt}\setlength{\parskip}{0pt}}



 


% ==============================================================================
% Preámbulo Institucional de Cooperación Técnica (Estilo GIZ - Corporate Design)
% Aplicado a clase KOMA-Script scrreprt con LuaLaTeX
% ==============================================================================

\usepackage{fontspec}
\usepackage{xcolor}
\usepackage{geometry}
\usepackage{scrlayer-scrpage}
\usepackage{tcolorbox}
\usepackage{booktabs}
\usepackage{tabularx}
\usepackage{hyperref}

% --- Localización y Nombres Oficiales en Español ---
\AtBeginDocument{
  \renewcommand{\contentsname}{Contenido}
  \renewcommand{\figurename}{Figura}
  \renewcommand{\tablename}{Tabla}
  \renewcommand{\chaptername}{Capítulo}
}

% --- Geometría de Página Oficial ---
\geometry{
  letterpaper,
  top=28mm,
  bottom=28mm,
  left=30mm,
  right=25mm,
  headheight=18pt,
  footskip=14mm
}

% --- Paleta Corporativa GIZ (Corporate Design Portal) ---
\definecolor{GIZRed}{RGB}{204, 6, 50}        % Primario Institucional: #CC0632
\definecolor{GIZDarkGray}{RGB}{51, 51, 51}    % Texto cuerpo y títulos sobrios: #333333
\definecolor{GIZLightGray}{RGB}{246, 246, 246}% Fondos de cajas y tablas: #F6F6F6
\definecolor{GIZBorderGray}{RGB}{210, 210, 210}
\definecolor{GIZAccentBlue}{RGB}{0, 92, 143}  % Acento secundario para datos/enlaces

% --- Tipografía Limpia Institucional ---
\setmainfont{TeX Gyre Heros}[
  Scale=0.95,
  Ligatures=TeX
]
\setsansfont{TeX Gyre Heros}[
  Scale=0.95,
  Ligatures=TeX
]
\setmonofont{TeX Gyre Cursor}[
  Scale=0.90
]

% Color por defecto del texto
\color{GIZDarkGray}

% --- Estilo de Encabezados y Pies de Página (scrlayer-scrpage) ---
\clearpairofpagestyles
\addtokomafont{pageheadfoot}{\sffamily\footnotesize\color{GIZDarkGray}}
\addtokomafont{pagenumber}{\sffamily\normalsize\bfseries\color{GIZRed}}

\ihead{\textnormal{Intervención Integral: Barrio El Capulí}}
\ohead{\textit{GIZ / CIS II — Convocatoria Ciudades que Transforman}}
\ifoot{\textcolor{gray}{\footnotesize Universidad Nacional de Loja — Economía \& Arquitectura}}
\ofoot{\pagemark}

\KOMAoptions{headsepline=true, footsepline=true}
\setkomafont{headsepline}{\color{GIZRed}}
\setkomafont{footsepline}{\color{GIZBorderGray}}

% --- Estilo de Títulos y Secciones KOMA-Script ---
\addtokomafont{disposition}{\sffamily\bfseries\color{GIZDarkGray}}
\addtokomafont{chapter}{\Huge\color{GIZRed}}
\addtokomafont{section}{\Large\color{GIZRed}}
\addtokomafont{subsection}{\large\color{GIZDarkGray}}
\addtokomafont{subsubsection}{\normalsize\bfseries\color{GIZDarkGray}}

% --- Cajas de Énfasis Metodológico GIZ (tcolorbox) ---
\tcbuselibrary{skins,breakable}

\newtcolorbox{gizbox}[1][]{
  colback=GIZLightGray,
  colframe=GIZRed,
  fonttitle=\bfseries\sffamily\color{white},
  colbacktitle=GIZRed,
  enhanced,
  boxrule=1pt,
  arc=2mm,
  left=4mm,
  right=4mm,
  top=3mm,
  bottom=3mm,
  breakable,
  #1
}

\newtcolorbox{gizcallout}[2][]{
  colback=white,
  colframe=GIZRed,
  fonttitle=\bfseries\sffamily\color{GIZRed},
  colbacktitle=white,
  enhanced,
  boxrule=0.8pt,
  leftrule=4pt,
  arc=0mm,
  title={#2},
  left=4mm,
  right=4mm,
  top=2mm,
  bottom=2mm,
  breakable,
  #1
}

% --- Ajuste de Hipervínculos ---
\hypersetup{
  colorlinks=true,
  linkcolor=GIZRed,
  filecolor=GIZRed,
  urlcolor=GIZAccentBlue,
  citecolor=GIZRed
}

\KOMAoption{captions}{tableheading}
\makeatletter
\@ifpackageloaded{caption}{}{\usepackage{caption}}
\AtBeginDocument{%
\ifdefined\contentsname
  \renewcommand*\contentsname{Table of contents}
\else
  \newcommand\contentsname{Table of contents}
\fi
\ifdefined\listfigurename
  \renewcommand*\listfigurename{List of Figures}
\else
  \newcommand\listfigurename{List of Figures}
\fi
\ifdefined\listtablename
  \renewcommand*\listtablename{List of Tables}
\else
  \newcommand\listtablename{List of Tables}
\fi
\ifdefined\figurename
  \renewcommand*\figurename{Figure}
\else
  \newcommand\figurename{Figure}
\fi
\ifdefined\tablename
  \renewcommand*\tablename{Table}
\else
  \newcommand\tablename{Table}
\fi
}
\@ifpackageloaded{float}{}{\usepackage{float}}
\floatstyle{ruled}
\@ifundefined{c@chapter}{\newfloat{codelisting}{h}{lop}}{\newfloat{codelisting}{h}{lop}[chapter]}
\floatname{codelisting}{Listing}
\newcommand*\listoflistings{\listof{codelisting}{List of Listings}}
\makeatother
\makeatletter
\makeatother
\makeatletter
\@ifpackageloaded{caption}{}{\usepackage{caption}}
\@ifpackageloaded{subcaption}{}{\usepackage{subcaption}}
\makeatother
\usepackage{bookmark}
\IfFileExists{xurl.sty}{\usepackage{xurl}}{} % add URL line breaks if available
\urlstyle{same}
% fallback for those not using the hyperref driver hyperxmp:
\makeatletter
\@ifundefined{xmpquote}{\newcommand{\xmpquote}[1]{#1}}{}
\makeatother
\hypersetup{
  pdftitle={PLANTILLA TÉCNICA INSTITUCIONAL: PERFIL DE PROYECTO URBANO SOSTENIBLE},
  pdfauthor={Equipo Técnico Interdisciplinario (Economía \& Arquitectura - UNL)},
  colorlinks=true,
  linkcolor={blue},
  filecolor={Maroon},
  citecolor={Blue},
  urlcolor={Blue},
  pdfcreator={LaTeX via pandoc}}


\title{PLANTILLA TÉCNICA INSTITUCIONAL: PERFIL DE PROYECTO URBANO
SOSTENIBLE}
\usepackage{etoolbox}
\makeatletter
\providecommand{\subtitle}[1]{% add subtitle to \maketitle
  \apptocmd{\@title}{\par {\large #1 \par}}{}{}
}
\makeatother
\subtitle{Estructura Operativa y Guía de Llenado bajo Estándar GIZ
(Ciudades Intermedias Sostenibles CIS II / Aceleración de Proyectos
Urbanos)}
\author{Equipo Técnico Interdisciplinario (Economía \& Arquitectura -
UNL)}
\date{2026-10-23}
\begin{document}
\maketitle

\renewcommand*\contentsname{Table of contents}
{
\setcounter{tocdepth}{2}
\tableofcontents
}

\chapter{PRESENTACIÓN Y GUÍA DE USO DE LA
PLANTILLA}\label{presentaciuxf3n-y-guuxeda-de-uso-de-la-plantilla}

\begin{quote}
\textbf{Objetivo de este documento:}\\
Esta plantilla estandarizada ha sido diseñada conforme a las directrices
de la \textbf{Cooperación Alemana al Desarrollo (GIZ)} ---en particular
los programas \emph{Ciudades Intermedias Sostenibles (CIS / CIS II)}, la
\emph{Guía para Elaborar Perfiles de Proyecto de Cambio Climático en
Ciudades (GIZ Ecuador, 2024)} y el marco \emph{FELICITY / TUMi}.

\textbf{Instrucciones para el equipo técnico:}\\
Los campos delimitados entre corchetes sombreados
\texttt{{[}Completar:\ ...{]}} representan las secciones editables donde
se debe vaciar y consolidar la información cuantitativa, espacial y
social recopilada en campo (KoBoToolbox, Censo CPV 2022 REDATAM, aforos
viales, cartografía QGIS y anteproyecto SbN).
\end{quote}

\begin{center}\rule{0.5\linewidth}{0.5pt}\end{center}

\chapter{MÓDULO 1: FICHA GENERAL DE IDENTIFICACIÓN DEL PROYECTO (PROJECT
FACTSHEET)}\label{muxf3dulo-1-ficha-general-de-identificaciuxf3n-del-proyecto-project-factsheet}

{\def\LTcaptype{none} % do not increment counter
\begin{longtable}[]{@{}
  >{\raggedright\arraybackslash}p{(\linewidth - 2\tabcolsep) * \real{0.5000}}
  >{\raggedright\arraybackslash}p{(\linewidth - 2\tabcolsep) * \real{0.5000}}@{}}
\toprule\noalign{}
\begin{minipage}[b]{\linewidth}\raggedright
Campo Metodológico
\end{minipage} & \begin{minipage}[b]{\linewidth}\raggedright
Descripción / Registro Técnico
\end{minipage} \\
\midrule\noalign{}
\endhead
\bottomrule\noalign{}
\endlastfoot
\textbf{Nombre del Proyecto:} &
\texttt{{[}Completar:\ Ej.\ Intervención\ Integral,\ Movilidad\ Activa\ y\ Resiliencia\ Hídrica\ en\ el\ Barrio\ El\ Capulí{]}} \\
\textbf{Localización Geográfica:} & Ecuador, Provincia de Loja, Cantón
Loja, Parroquia Punzara, Sector Capulí Loma. \\
\textbf{Coordenadas de Referencia (WGS84 / UTM 17S):} &
\texttt{{[}Completar:\ X:\ ...\ ,\ Y:\ ...\ /\ Polígono\ delimitado\ en\ QGIS{]}} \\
\textbf{Entidad Proponente / Promotora:} & Universidad Nacional de Loja
(UNL) en articulación con el GAD Municipal de Loja y la Directiva
Barrial. \\
\textbf{Área de Influencia Territorial:} &
\texttt{{[}Completar:\ Ej.\ Polígono\ de\ intervención\ directa:\ XX\ hectáreas\ /\ Polígono\ indirecto:\ XX\ ha{]}} \\
\textbf{Población Beneficiaria Directa:} &
\texttt{{[}Completar:\ N°\ de\ habitantes\ del\ Barrio\ El\ Capulí\ según\ Censo\ 2022\ y\ padrón\ local{]}} \\
\textbf{Población Beneficiaria Indirecta:} &
\texttt{{[}Completar:\ Población\ del\ área\ sur\ de\ Loja,\ comunidad\ universitaria\ La\ Argelia,\ usuarios\ de\ la\ Av.\ Reinaldo\ Espinosa{]}} \\
\textbf{Tipología de Intervención GIZ:} & ☐ Movilidad Activa y Calles
Completas ☐ Soluciones basadas en la Naturaleza (SbN / SUDS) ☐ Espacio
Público y Sistema Verde Urbano ☒ Proyecto Integral / Multicomponente \\
\textbf{Horizonte Temporal del Proyecto:} & Fase de formulación: 2026 \\
\end{longtable}
}

\begin{center}\rule{0.5\linewidth}{0.5pt}\end{center}

\chapter{MÓDULO 2: DIAGNÓSTICO TERRITORIAL Y LÍNEA BASE
CUANTITATIVA}\label{muxf3dulo-2-diagnuxf3stico-territorial-y-luxednea-base-cuantitativa}

\section{Delimitación Espacial y Contexto
Geográfico}\label{delimitaciuxf3n-espacial-y-contexto-geogruxe1fico}

\emph{Describa las características fisiográficas, topográficas y los
límites del sector de intervención.}

\texttt{{[}Completar\ texto:\ Describir\ la\ localización\ en\ la\ parroquia\ Punzara,\ la\ topografía\ accidentada\ de\ Capulí\ Loma\ con\ pendientes\ predominantes\ entre\ 15\%\ y\ 35\%,\ y\ su\ condición\ de\ borde\ urbano\ periurbano.{]}}

\section{Caracterización Socioeconómica y Déficit Habitacional
(Microdatos Censales /
KoBoToolbox)}\label{caracterizaciuxf3n-socioeconuxf3mica-y-duxe9ficit-habitacional-microdatos-censales-kobotoolbox}

\emph{Inserte la línea base cuantitativa obtenida del procesamiento del
Censo CPV 2022 (REDATAM) y las encuestas de campo.}

{\def\LTcaptype{none} % do not increment counter
\begin{longtable}[]{@{}
  >{\raggedright\arraybackslash}p{(\linewidth - 6\tabcolsep) * \real{0.2500}}
  >{\raggedright\arraybackslash}p{(\linewidth - 6\tabcolsep) * \real{0.2500}}
  >{\raggedright\arraybackslash}p{(\linewidth - 6\tabcolsep) * \real{0.2500}}
  >{\raggedright\arraybackslash}p{(\linewidth - 6\tabcolsep) * \real{0.2500}}@{}}
\toprule\noalign{}
\begin{minipage}[b]{\linewidth}\raggedright
Variable / Indicador
\end{minipage} & \begin{minipage}[b]{\linewidth}\raggedright
Valor Levantado (Línea Base)
\end{minipage} & \begin{minipage}[b]{\linewidth}\raggedright
Comparativa Cantonal / Parroquial
\end{minipage} & \begin{minipage}[b]{\linewidth}\raggedright
Fuente de Información
\end{minipage} \\
\midrule\noalign{}
\endhead
\bottomrule\noalign{}
\endlastfoot
\textbf{Hogares con Necesidades Básicas Insatisfechas (NBI):} &
\texttt{{[}Completar:\ ej.\ 38.4\%{]}} &
\texttt{{[}Completar:\ Cantón\ Loja:\ XX\%{]}} & Censo INEC 2022
(REDATAM) \\
\textbf{Cobertura de Red Pública de Agua Potable:} &
\texttt{{[}Completar:\ ej.\ Déficit\ \textgreater{}\ 80\%\ /\ intermitente{]}}
& \texttt{{[}Completar:\ Parroquia:\ XX\%{]}} & KoBoToolbox / Censo
2022 \\
\textbf{Acceso a Alcantarillado Sanitario / Pluvial:} &
\texttt{{[}Completar:\ Cobertura\ XX\%\ /\ Pozos\ sépticos{]}} &
\texttt{{[}Completar:\ Cantón:\ XX\%{]}} & Levantamiento de campo \\
\textbf{Tipología y Calidad de la Vivienda:} &
\texttt{{[}Completar:\ \%\ vivienda\ precaria\ /\ ladrillo\ /\ adobe{]}}
& \texttt{{[}Completar{]}} & Fichas de Habitabilidad KoBo \\
\textbf{Estructura Demográfica (Edad / Género):} &
\texttt{{[}Completar:\ \%\ infancia,\ adultos\ mayores,\ género{]}} &
\texttt{{[}Completar{]}} & Censo INEC 2022 \\
\end{longtable}
}

\section{Diagnóstico del Sistema de Conectividad y Movilidad
Urbana}\label{diagnuxf3stico-del-sistema-de-conectividad-y-movilidad-urbana}

\emph{Analice las condiciones de transitabilidad, transporte público y
seguridad vial.}

\begin{enumerate}
\def\labelenumi{\arabic{enumi}.}
\tightlist
\item
  \textbf{Jerarquización Vial y Estado de Calzada:}\\
  \texttt{{[}Completar:\ Análisis\ de\ la\ Av.\ Reinaldo\ Espinosa\ como\ conector\ artificial\ primario\ y\ red\ de\ calles\ terciarias/pasajes\ en\ tierra\ o\ lastre\ sin\ aceras.{]}}
\item
  \textbf{Aforos y Modos de Desplazamiento:}\\
  \texttt{{[}Completar:\ Resultados\ de\ los\ aforos\ vehiculares\ y\ peatonales,\ frecuencias\ de\ buses\ urbanos\ y\ modos\ de\ transporte\ utilizados\ por\ los\ moradores.{]}}
\item
  \textbf{Isócronas de Caminabilidad y Accesibilidad:}\\
  \texttt{{[}Completar:\ Tiempos\ de\ desplazamiento\ a\ pie\ hacia\ el\ Campus\ Universitario\ La\ Argelia,\ paradas\ de\ bus\ más\ cercanas\ y\ equipamientos\ básicos\ (10-15\ min).{]}}
\end{enumerate}

\section{Vulnerabilidad Física, Pendientes y Riesgo
Hídrico}\label{vulnerabilidad-fuxedsica-pendientes-y-riesgo-huxeddrico}

\emph{Detalle los riesgos por escorrentía pluvial y estabilidad de
laderas.}

\texttt{{[}Completar:\ Detallar\ los\ puntos\ críticos\ de\ acumulación\ de\ escorrentías\ pluviales\ en\ época\ invernal,\ erosión\ de\ calzadas\ no\ pavimentadas,\ falta\ de\ subdrenajes\ y\ desbordamientos\ hacia\ la\ microcuenca\ del\ río\ Capulí.{]}}

\begin{center}\rule{0.5\linewidth}{0.5pt}\end{center}

\chapter{MÓDULO 3: TEORÍA DE CAMBIO Y ESTRUCTURA LÓGICA (METODOLOGÍA
GIZ)}\label{muxf3dulo-3-teoruxeda-de-cambio-y-estructura-luxf3gica-metodologuxeda-giz}

\section{Árbol de Problemas (Causa -
Efecto)}\label{uxe1rbol-de-problemas-causa---efecto}

\begin{itemize}
\tightlist
\item
  \textbf{Efecto Final:} Deterioro de la calidad de vida, vulnerabilidad
  socioambiental y fragmentación socio-espacial de los habitantes del
  Barrio El Capulí.

  \begin{itemize}
  \tightlist
  \item
    \textbf{Efecto 1.1:} Altos costos y tiempos de desplazamiento;
    aislamiento de grupos vulnerables.
  \item
    \textbf{Efecto 1.2:} Inseguridad vial para peatones y ciclistas;
    alta siniestralidad potencial.
  \item
    \textbf{Efecto 1.3:} Pérdida de biodiversidad urbana, contaminación
    hídrica y anegamiento constante.
  \end{itemize}
\item
  \textbf{PROBLEMA CENTRAL:} \textbf{Deficiente conectividad vial,
  precariedad del espacio público y alta vulnerabilidad ante amenazas
  hídricas en el Barrio El Capulí.}

  \begin{itemize}
  \tightlist
  \item
    \textbf{Causa Directa 1:} Infraestructura vial incompleta,
    desarticulada y carente de criterios de movilidad activa 8-80.

    \begin{itemize}
    \tightlist
    \item
      \emph{Causa Raíz 1.1:} Crecimiento urbano no planificado en ladera
      sin dotación reglamentaria de aceras.
    \end{itemize}
  \item
    \textbf{Causa Directa 2:} Inexistencia de un sistema de drenaje
    pluvial sostenible y desprotección de la microcuenca.

    \begin{itemize}
    \tightlist
    \item
      \emph{Causa Raíz 2.1:} Ausencia de Soluciones basadas en la
      Naturaleza (SbN / SUDS) para regular escorrentías.
    \end{itemize}
  \item
    \textbf{Causa Directa 3:} Carencia de equipamiento público
    estructurante e integración ambiental.

    \begin{itemize}
    \tightlist
    \item
      \emph{Causa Raíz 3.1:} Falta de consolidación del Parque de río
      Capulí dentro del Sistema Verde Urbano cantonal.
    \end{itemize}
  \end{itemize}
\end{itemize}

\section{Árbol de Objetivos (Medios -
Fines)}\label{uxe1rbol-de-objetivos-medios---fines}

\begin{itemize}
\tightlist
\item
  \textbf{Fin Último:} Incremento del bienestar, integración territorial
  sostenible y resiliencia climática en el Barrio El Capulí.
\item
  \textbf{OBJETIVO GENERAL:} \textbf{Implementar un modelo de
  intervención integral urbana y ambiental en El Capulí mediante
  conectividad activa, soluciones basadas en la naturaleza y
  consolidación de espacio público resiliente.}

  \begin{itemize}
  \tightlist
  \item
    \textbf{Componente 1:} Red de movilidad activa segura, continua e
    inclusiva (Calle Completa).
  \item
    \textbf{Componente 2:} Sistema de amortiguamiento y drenaje pluvial
    sostenible en laderas (SbN).
  \item
    \textbf{Componente 3:} Parque de río Capulí consolidado como
    infraestructura verde comunitaria activa.
  \end{itemize}
\end{itemize}

\section{Alineación Estratégica con Instrumentos de
Planificación}\label{alineaciuxf3n-estratuxe9gica-con-instrumentos-de-planificaciuxf3n}

{\def\LTcaptype{none} % do not increment counter
\begin{longtable}[]{@{}
  >{\raggedright\arraybackslash}p{(\linewidth - 4\tabcolsep) * \real{0.3333}}
  >{\raggedright\arraybackslash}p{(\linewidth - 4\tabcolsep) * \real{0.3333}}
  >{\raggedright\arraybackslash}p{(\linewidth - 4\tabcolsep) * \real{0.3333}}@{}}
\toprule\noalign{}
\begin{minipage}[b]{\linewidth}\raggedright
Instrumento / Marco de Política
\end{minipage} & \begin{minipage}[b]{\linewidth}\raggedright
Eje / Meta Específica
\end{minipage} & \begin{minipage}[b]{\linewidth}\raggedright
Alineación del Proyecto
\end{minipage} \\
\midrule\noalign{}
\endhead
\bottomrule\noalign{}
\endlastfoot
\textbf{Agenda Hábitat Sostenible del Ecuador (AHSE 2036):} & Eje
Ciudades Resilientes y Accesibles & Intervención integral en periferia
urbana y SbN. \\
\textbf{Plan de Desarrollo y Ordenamiento Territorial (PDOT Loja):} &
Eje Territorial y Movilidad Sostenible & Jerarquización vial y reducción
de brechas periféricas. \\
\textbf{Plan del Sistema Verde Urbano de Loja:} & Proyecto Estructurante
N° 11 & Recuperación y anteproyecto del Parque de río Capulí. \\
\textbf{Objetivos de Desarrollo Sostenible (ODS):} & ODS 11 (Ciudades
Sostenibles), ODS 13 (Acción por el Clima), ODS 6 (Agua Limpia) &
Mitigación climática y acceso a infraestructura básica. \\
\end{longtable}
}

\begin{center}\rule{0.5\linewidth}{0.5pt}\end{center}

\chapter{MÓDULO 4: PROPUESTA TÉCNICA Y COMPONENTES DE
INTERVENCIÓN}\label{muxf3dulo-4-propuesta-tuxe9cnica-y-componentes-de-intervenciuxf3n}

\section{Componente 1: Conectividad Vial y Movilidad Activa (Calle
Completa
8-80)}\label{componente-1-conectividad-vial-y-movilidad-activa-calle-completa-8-80}

\emph{Describa el diseño geométrico, perfiles viales y distribución
modal propuestos.}

\begin{itemize}
\tightlist
\item
  \textbf{Perfil Transversal Tipo (Calle Completa):}\\
  \texttt{{[}Completar:\ Ancho\ de\ calzada\ vehicular,\ ancho\ de\ veredas\ peatonales\ accesibles\ con\ guía\ podotáctil,\ ciclovía\ segregada/compartida,\ arbolado\ viario\ y\ franjas\ de\ amortiguamiento.{]}}
\item
  \textbf{Conexión Intermodal hacia el Campus La Argelia:}\\
  \texttt{{[}Completar:\ Trazado\ de\ senderos\ de\ ladera\ escalonados\ y\ rampas\ universales\ que\ reduzcan\ el\ tiempo\ de\ traslado\ a\ pie\ de\ 25\ min\ a\ 10\ min.{]}}
\item
  \textbf{Seguridad Vial y Pacificación de Tráfico:}\\
  \texttt{{[}Completar:\ Reductores\ de\ velocidad\ tipo\ lomo\ de\ buey/plataforma\ única\ en\ cruces\ peatonales,\ señalética\ horizontal/vertical\ e\ iluminación\ LED\ solar.{]}}
\end{itemize}

\section{Componente 2: Resiliencia Hídrica y Soluciones basadas en la
Naturaleza (SbN /
SUDS)}\label{componente-2-resiliencia-huxeddrica-y-soluciones-basadas-en-la-naturaleza-sbn-suds}

\emph{Detalle las tecnologías ambientales propuestas para el manejo de
aguas de lluvia y laderas.}

\begin{itemize}
\tightlist
\item
  \textbf{Jardines de Lluvia y Cunetas Verdes Biofiltrantes:}\\
  \texttt{{[}Completar:\ Ubicación\ en\ márgenes\ de\ vía\ para\ retener\ el\ primer\ flujo\ de\ escorrentía\ pluvial\ y\ filtrar\ sedimentos.{]}}
\item
  \textbf{Pavimentos Permeables y Zanjas de Infiltración:}\\
  \texttt{{[}Completar:\ Materiales\ propuestos\ en\ aceras\ y\ senderos\ para\ evitar\ la\ escorrentía\ superficial\ acelerada.{]}}
\item
  \textbf{Bioingeniería de Estabilización de Taludes:}\\
  \texttt{{[}Completar:\ Viveros\ de\ especies\ nativas\ (vetiver,\ aliso,\ sigse)\ para\ protección\ de\ taludes\ con\ pendientes\ superiores\ al\ 25\%.{]}}
\end{itemize}

\section{Componente 3: Espacio Público y Parque Fluvial de río
Capulí}\label{componente-3-espacio-puxfablico-y-parque-fluvial-de-ruxedo-capuluxed}

\emph{Describa el anteproyecto paisajístico y comunitario.}

\begin{itemize}
\tightlist
\item
  \textbf{Zonificación del Parque Fluvial:}\\
  \texttt{{[}Completar:\ Zona\ de\ conservación\ ribereña,\ senderos\ ecológicos,\ áreas\ lúdico-deportivas,\ mirador\ paisajístico\ y\ plazas\ de\ encuentro\ barrial.{]}}
\item
  \textbf{Equipamiento Urbano Sostenible:}\\
  \texttt{{[}Completar:\ Mobiliario\ urbano\ en\ madera\ plástica\ reciclada,\ iluminación\ fotovoltaica,\ puntos\ limpios\ y\ señalética\ didáctica\ ambiental.{]}}
\end{itemize}

\begin{center}\rule{0.5\linewidth}{0.5pt}\end{center}

\chapter{MÓDULO 5: MODELO DE GESTIÓN, PRESUPUESTO Y SOSTENIBILIDAD
(CAPEX /
OPEX)}\label{muxf3dulo-5-modelo-de-gestiuxf3n-presupuesto-y-sostenibilidad-capex-opex}

\section{Estimación de Costos de Inversión Inicial
(CAPEX)}\label{estimaciuxf3n-de-costos-de-inversiuxf3n-inicial-capex}

{\def\LTcaptype{none} % do not increment counter
\begin{longtable}[]{@{}
  >{\centering\arraybackslash}p{(\linewidth - 10\tabcolsep) * \real{0.1724}}
  >{\raggedright\arraybackslash}p{(\linewidth - 10\tabcolsep) * \real{0.1379}}
  >{\centering\arraybackslash}p{(\linewidth - 10\tabcolsep) * \real{0.1724}}
  >{\centering\arraybackslash}p{(\linewidth - 10\tabcolsep) * \real{0.1724}}
  >{\centering\arraybackslash}p{(\linewidth - 10\tabcolsep) * \real{0.1724}}
  >{\centering\arraybackslash}p{(\linewidth - 10\tabcolsep) * \real{0.1724}}@{}}
\toprule\noalign{}
\begin{minipage}[b]{\linewidth}\centering
N°
\end{minipage} & \begin{minipage}[b]{\linewidth}\raggedright
Componente / Rubro General
\end{minipage} & \begin{minipage}[b]{\linewidth}\centering
Unidad
\end{minipage} & \begin{minipage}[b]{\linewidth}\centering
Cantidad
\end{minipage} & \begin{minipage}[b]{\linewidth}\centering
Costo Unitario Est. (USD)
\end{minipage} & \begin{minipage}[b]{\linewidth}\centering
Costo Total (USD)
\end{minipage} \\
\midrule\noalign{}
\endhead
\bottomrule\noalign{}
\endlastfoot
\textbf{1} & Estudios definitivos y geotecnia & Glb & 1 &
\texttt{{[}Completar{]}} & \texttt{{[}Completar{]}} \\
\textbf{2} & Movilidad Activa y Veredas Inclusivas & m² / km &
\texttt{{[}Completar{]}} & \texttt{{[}Completar{]}} &
\texttt{{[}Completar{]}} \\
\textbf{3} & Soluciones basadas en la Naturaleza (SUDS) & m² &
\texttt{{[}Completar{]}} & \texttt{{[}Completar{]}} &
\texttt{{[}Completar{]}} \\
\textbf{4} & Infraestructura Verde y Parque de río Capulí & ha &
\texttt{{[}Completar{]}} & \texttt{{[}Completar{]}} &
\texttt{{[}Completar{]}} \\
\textbf{5} & Fiscalización y Control Ambiental (5\%) & Glb & 1 &
\texttt{{[}Completar{]}} & \texttt{{[}Completar{]}} \\
\textbf{--} & \textbf{TOTAL ESTIMADO DE INVERSIÓN (CAPEX):} & -- & -- &
-- & \textbf{\texttt{{[}Completar\ USD{]}}} \\
\end{longtable}
}

\section{Costos de Operación y Mantenimiento a Largo Plazo
(OPEX)}\label{costos-de-operaciuxf3n-y-mantenimiento-a-largo-plazo-opex}

{\def\LTcaptype{none} % do not increment counter
\begin{longtable}[]{@{}
  >{\raggedright\arraybackslash}p{(\linewidth - 6\tabcolsep) * \real{0.2500}}
  >{\raggedright\arraybackslash}p{(\linewidth - 6\tabcolsep) * \real{0.2500}}
  >{\raggedright\arraybackslash}p{(\linewidth - 6\tabcolsep) * \real{0.2500}}
  >{\raggedright\arraybackslash}p{(\linewidth - 6\tabcolsep) * \real{0.2500}}@{}}
\toprule\noalign{}
\begin{minipage}[b]{\linewidth}\raggedright
Periodo / Horizonte
\end{minipage} & \begin{minipage}[b]{\linewidth}\raggedright
Actividades Clave de Mantenimiento
\end{minipage} & \begin{minipage}[b]{\linewidth}\raggedright
Costo Anual Promedio (USD)
\end{minipage} & \begin{minipage}[b]{\linewidth}\raggedright
Responsable Operativo
\end{minipage} \\
\midrule\noalign{}
\endhead
\bottomrule\noalign{}
\endlastfoot
\textbf{Corto Plazo (Años 1 al 5):} & Limpieza de sedimentos en SUDS,
podas de arbolado, reposición de pintura vial. &
\texttt{{[}Completar\ USD{]}} & GAD Loja / Comité Barrial \\
\textbf{Mediano Plazo (Años 6 al 10):} & Mantenimiento de carpetas y
juntas en senderos, reforestación complementaria. &
\texttt{{[}Completar\ USD{]}} & GAD Municipal de Loja \\
\textbf{Largo Plazo (Años 11 al 20):} & Rehabilitación mayor de
mobiliario y obras civiles de contención. &
\texttt{{[}Completar\ USD{]}} & Presupuesto Participativo GAD \\
\end{longtable}
}

\section{Esquema Institucional y Gobernanza del
Proyecto}\label{esquema-institucional-y-gobernanza-del-proyecto}

\begin{itemize}
\tightlist
\item
  \textbf{Nivel Directivo / Financiero:} GAD Municipal de Loja y líneas
  de financiamiento verde (BDE, Fondos Climáticos, GIZ).
\item
  \textbf{Nivel Técnico y Asistencia Académica:} Universidad Nacional de
  Loja (carreras de Economía y Arquitectura - Vinculación con la
  Sociedad).
\item
  \textbf{Nivel de Custodia Comunitaria:} Junta Directiva del Barrio El
  Capulí (Comité Barrial de Vigilancia y Cuidado del Parque).
\end{itemize}

\begin{center}\rule{0.5\linewidth}{0.5pt}\end{center}

\chapter{MÓDULO 6: MATRIZ DE MARCO LÓGICO, INDICADORES Y MONITOREO
(MRV)}\label{muxf3dulo-6-matriz-de-marco-luxf3gico-indicadores-y-monitoreo-mrv}

{\def\LTcaptype{none} % do not increment counter
\begin{longtable}[]{@{}
  >{\raggedright\arraybackslash}p{(\linewidth - 10\tabcolsep) * \real{0.1667}}
  >{\raggedright\arraybackslash}p{(\linewidth - 10\tabcolsep) * \real{0.1667}}
  >{\raggedright\arraybackslash}p{(\linewidth - 10\tabcolsep) * \real{0.1667}}
  >{\raggedright\arraybackslash}p{(\linewidth - 10\tabcolsep) * \real{0.1667}}
  >{\raggedright\arraybackslash}p{(\linewidth - 10\tabcolsep) * \real{0.1667}}
  >{\raggedright\arraybackslash}p{(\linewidth - 10\tabcolsep) * \real{0.1667}}@{}}
\toprule\noalign{}
\begin{minipage}[b]{\linewidth}\raggedright
Nivel de Objetivos
\end{minipage} & \begin{minipage}[b]{\linewidth}\raggedright
Indicador Cuantitativo
\end{minipage} & \begin{minipage}[b]{\linewidth}\raggedright
Línea Base (2026)
\end{minipage} & \begin{minipage}[b]{\linewidth}\raggedright
Meta del Proyecto
\end{minipage} & \begin{minipage}[b]{\linewidth}\raggedright
Fuente de Verificación
\end{minipage} & \begin{minipage}[b]{\linewidth}\raggedright
Supuestos / Riesgos
\end{minipage} \\
\midrule\noalign{}
\endhead
\bottomrule\noalign{}
\endlastfoot
\textbf{Fin:} & Reducción de la vulnerabilidad socio-urbana & NBI:
38.4\% & Reducción en 15\% & Censo / Encuestas socioeconómicas &
Estabilidad institucional y macroeconómica \\
\textbf{Propósito:} & Población con acceso a movilidad segura y espacio
público de calidad & 0 m² de parque consolidado / 0\% veredas accesibles
& 100\% de cobertura en eje central; 1 Parque Fluvial habilitado &
Auditoría urbana y actas de entrega-recepción & Aprobación del PDOT y
asignación presupuestaria \\
\textbf{Componente 1:} & Kilómetros de senderos y veredas 8-80
implementadas & 0.0 km con diseño universal &
\texttt{{[}Completar:\ ej.\ 2.5\ km{]}} & Planos As-Built y
fiscalización de obra & Cumplimiento contractual de contratistas \\
\textbf{Componente 2:} & Caudal pico de escorrentía amortiguado por SbN
& 0\% retención & Reducción de al menos 30\% del volumen pico pluvial &
Medición hidrológica de campo & Fenómenos climáticos extremos
atípicos \\
\textbf{Componente 3:} & Superficie de áreas verdes ribereñas
recuperadas & Microcuenca degradada &
\texttt{{[}Completar:\ ej.\ 1.8\ ha{]}} & Levantamiento topográfico y
ortofotos satelitales & Respaldo y apropiación de los vecinos \\
\end{longtable}
}

\begin{center}\rule{0.5\linewidth}{0.5pt}\end{center}

\chapter{MÓDULO 7: ANEXOS TÉCNICOS OBLIGATORIOS
REQUERIDOS}\label{muxf3dulo-7-anexos-tuxe9cnicos-obligatorios-requeridos}

\begin{itemize}
\tightlist
\item
  \textbf{Anexo A:} Cartografía Temática en Alta Resolución (Planos
  QGIS: Pendientes, Usos de Suelo, Isócronas y Propuesta).
\item
  \textbf{Anexo B:} Base de Datos Tabulada de Encuestas KoBoToolbox
  (Microdatos de campo anonimizados).
\item
  \textbf{Anexo C:} Tabulados de Microdatos Censales REDATAM (INEC CPV
  2022).
\item
  \textbf{Anexo D:} Memorias de Cálculo Hidrológico y Dimensionamiento
  Preliminar de SUDS.
\item
  \textbf{Anexo E:} Planos Arquitectónicos y Secciones Viales del
  Anteproyecto 8-80.
\item
  \textbf{Anexo F:} Registro Fotográfico Georreferenciado y Actas de
  Talleres Comunitarios Participativos.
\end{itemize}




\end{document}
